\documentclass[10pt,a4paper]{amsart}
\usepackage{amsmath,amssymb,amsthm,graphicx}
\usepackage{verbatim}
\newenvironment{reference}{\comment}{\endcomment}
%\newenvironment{reference}{}{}
\newenvironment{slogan}{\comment}{\endcomment}
\newenvironment{history}{\comment}{\endcomment}

% For commutative diagrams we use Xy-pic
\usepackage[all]{xy}

% We use 2cell for 2-commutative diagrams.
\xyoption{2cell}
% Theorem environments.
%
\theoremstyle{plain}
\newtheorem{theorem}[subsection]{Theorem}
\newtheorem{proposition}[subsection]{Proposition}
\newtheorem{lemma}[subsection]{Lemma}

\theoremstyle{definition}
\newtheorem{definition}[subsection]{Definition}
\newtheorem{example}[subsection]{Example}
\newtheorem{exercise}[subsection]{Exercise}
\newtheorem{situation}[subsection]{Situation}

\theoremstyle{remark}
\newtheorem{remark}[subsection]{Remark}
\newtheorem{remarks}[subsection]{Remarks}

\numberwithin{equation}{subsection}

% Macros
%
\def\lim{\mathop{\mathrm{lim}}\nolimits}
\def\colim{\mathop{\mathrm{colim}}\nolimits}
\def\Spec{\mathop{\mathrm{Spec}}}
\def\Hom{\mathop{\mathrm{Hom}}\nolimits}
\def\Ext{\mathop{\mathrm{Ext}}\nolimits}
\def\SheafHom{\mathop{\mathcal{H}\!\mathit{om}}\nolimits}
\def\SheafExt{\mathop{\mathcal{E}\!\mathit{xt}}\nolimits}
\def\Sch{\mathit{Sch}}
\def\Mor{\mathop{\mathrm{Mor}}\nolimits}
\def\Ob{\mathop{\mathrm{Ob}}\nolimits}
\def\Sh{\mathop{\mathit{Sh}}\nolimits}
\def\NL{\mathop{N\!L}\nolimits}
\def\CH{\mathop{\mathrm{CH}}\nolimits}
\def\proetale{{pro\text{-}\acute{e}tale}}
\def\etale{{\acute{e}tale}}
\def\QCoh{\mathit{QCoh}}
\def\Ker{\mathop{\mathrm{Ker}}}
\def\Im{\mathop{\mathrm{Im}}}
\def\Coker{\mathop{\mathrm{Coker}}}
\def\Coim{\mathop{\mathrm{Coim}}}

% Boxtimes
%
\DeclareMathSymbol{\boxtimes}{\mathbin}{AMSa}{"02}

%
% Macros for moduli stacks/spaces
%
\def\QCohstack{\mathcal{QC}\!\mathit{oh}}
\def\Cohstack{\mathcal{C}\!\mathit{oh}}
\def\Spacesstack{\mathcal{S}\!\mathit{paces}}
\def\Quotfunctor{\mathrm{Quot}}
\def\Hilbfunctor{\mathrm{Hilb}}
\def\Curvesstack{\mathcal{C}\!\mathit{urves}}
\def\Polarizedstack{\mathcal{P}\!\mathit{olarized}}
\def\Complexesstack{\mathcal{C}\!\mathit{omplexes}}
% \Pic is the operator that assigns to X its picard group, usage \Pic(X)
% \Picardstack_{X/B} denotes the Picard stack of X over B
% \Picardfunctor_{X/B} denotes the Picard functor of X over B
\def\Pic{\mathop{\mathrm{Pic}}\nolimits}
\def\Picardstack{\mathcal{P}\!\mathit{ic}}
\def\Picardfunctor{\mathrm{Pic}}
\def\Deformationcategory{\mathcal{D}\!\mathit{ef}}
\usepackage{cleveref}
\crefname{originallemma}{original lemma}{original lemmas}
\Crefname{originallemma}{Original lemma}{Original lemmas}
\makeatletter
\newlabel{lemma-localize-cocontinuous}{{03EG}{source}{}{sourcepointer-1}{}}
\newlabel{lemma-localize-cocontinuous@cref}{{[originallemma][0][]03EG}{[1][0][]source}{}{}{}}
\makeatother
\crefname{originallemma}{original lemma}{original lemmas}
\Crefname{originallemma}{Original lemma}{Original lemmas}
\makeatletter
\newlabel{lemma-equivalence}{{03A0}{source}{}{sourcepointer-2}{}}
\newlabel{lemma-equivalence@cref}{{[originallemma][0][]03A0}{[1][0][]source}{}{}{}}
\makeatother
\title{Localizing a special cocontinuous functor stays special cocontinuous}
\author{The Stacks Project Authors}
\date{stacks-project:sites.tex:local-snapshot-20261012; GNU FDL 1.2 or later}
\begin{document}
\maketitle
\noindent\textit{Editorial note.} All delivered mathematical text is the author's own, copied from the pinned Stacks Project chapter file sites-cohomology.tex; no statement or proof step was rewritten, reordered or removed.\par\medskip
\noindent\textit{Editorial note.} The proof is the author's complete argument; every step the author writes out is inside the counted span and no closing sentence is a placeholder.\par\medskip
\noindent\textit{Editorial note.} Every labelled reference occurring inside the counted spans is carried in assembly.external\_references as a non-load-bearing navigation pointer (same-chapter ones at their definition site, cross-chapter ones at the invocation inside the counted span), so no load-bearing mathematics is replaced by a stub. No bibliography citation occurs inside the counted spans.\par\medskip
\noindent\textit{Editorial note.} The checked chapter list for same-name collisions is the set of stacks chapters staged in this repository: chow, cohomology, divisors, etale-cohomology, flat, more-algebra, more-morphisms, morphisms, spaces-more-morphisms, varieties together with this card's chapter sites-cohomology.\par\medskip
\noindent\textit{Editorial note.} The preamble is the chapter's own retained minimal preamble (sources/stacks-sites-cohomology/retained-preamble.tex), a verbatim excerpt of the native preamble: the Xy-pic setup and the native theorem declarations with the mathematical macros.\par\medskip
\noindent\textit{Allowed standard prerequisites (editorial).}\begin{itemize}
\item The five conditions characterizing an equivalence of topoi. Same chapter, Lemma lemma-equivalence. Supplies the five properties verified for the induced functor on the localized sites.
\item The localization of a cocontinuous functor and its commutativity. Same chapter, Lemma lemma-localize-cocontinuous. Supplies the existence of the induced functor and the commutativity of the displayed diagrams.
\item Continuity and cocontinuity. Same chapter, Definition definition-cocontinuous and the definition of a continuous functor. Supplies the behaviour of coverings and refinements under the functor.
\end{itemize}\medskip
\ifdefined\hypertarget\hypertarget{sourcepointer-1}{}\fi
\noindent\textit{Source reference.} Original lemma 03EG: sources/stacks-sites/pinned-original-source.tex, source lines 7129--7129. Navigation pointer invoked at proof line 7129; the counted span contains the pointer only and the cited result is not reproduced inside it, so the pointer is registered as a non-load-bearing navigation stub in the established form for Stacks material. No bibliography citation is involved.\par\medskip
\ifdefined\hypertarget\hypertarget{sourcepointer-2}{}\fi
\noindent\textit{Source reference.} Original lemma 03A0: sources/stacks-sites/pinned-original-source.tex, source lines 7108--7108. Navigation pointer invoked at proof line 7108; the counted span contains the pointer only and the cited result is not reproduced inside it, so the pointer is registered as a non-load-bearing navigation stub in the established form for Stacks material. No bibliography citation is involved.\par\medskip
\begin{lemma}
\label{lemma-localize-special-cocontinuous}
Let $\mathcal{C}$, $\mathcal{D}$ be sites.
Let $u : \mathcal{C} \to \mathcal{D}$ be a special cocontinuous functor.
For every object $U$ of $\mathcal{C}$ we have a commutative diagram
$$
\xymatrix{
\mathcal{C}/U \ar[r]_{j_U} \ar[d] & \mathcal{C} \ar[d]^u \\
\mathcal{D}/u(U) \ar[r]^-{j_{u(U)}} & \mathcal{D}
}
$$
as in Lemma \ref{lemma-localize-cocontinuous}.
The left vertical arrow is a special cocontinuous functor.
Hence in the commutative diagram of topoi
$$
\xymatrix{
\Sh(\mathcal{C}/U) \ar[r]_{j_U} \ar[d] &
\Sh(\mathcal{C}) \ar[d]^u \\
\Sh(\mathcal{D}/u(U)) \ar[r]^-{j_{u(U)}} &
\Sh(\mathcal{D})
}
$$
the vertical arrows are equivalences.
\end{lemma}

\begin{proof}
We have seen the existence and commutativity of the diagrams in
Lemma \ref{lemma-localize-cocontinuous}. We have to check
hypotheses (1) -- (5) of Lemma \ref{lemma-equivalence} for the
induced functor $u : \mathcal{C}/U \to \mathcal{D}/u(U)$.
This is completely mechanical.

\medskip\noindent
Property (1). This is Lemma \ref{lemma-localize-cocontinuous}.

\medskip\noindent
Property (2). Let $\{U_i'/U \to U'/U\}_{i \in I}$ be a covering
of $U'/U$ in $\mathcal{C}/U$. Because $u$ is continuous we see that
$\{u(U_i')/u(U) \to u(U')/u(U)\}_{i \in I}$ is a covering
of $u(U')/u(U)$ in $\mathcal{D}/u(U)$. Hence (2) holds
for $u : \mathcal{C}/U \to \mathcal{D}/u(U)$.

\medskip\noindent
Property (3). Let $a, b : U''/U \to U'/U$ in $\mathcal{C}/U$
be morphisms such that $u(a) = u(b)$ in $\mathcal{D}/u(U)$.
Because $u$ satisfies (3) we see there exists a covering
$\{f_i : U''_i \to U''\}$ in $\mathcal{C}$ such that
$a \circ f_i = b \circ f_i$. This gives a covering
$\{f_i : U''_i/U \to U''/U\}$ in $\mathcal{C}/U$ such that
$a \circ f_i = b \circ f_i$. Hence (3) holds
for $u : \mathcal{C}/U \to \mathcal{D}/u(U)$.

\medskip\noindent
Property (4). Let $U''/U, U'/U \in \Ob(\mathcal{C}/U)$ and
a morphism $c : u(U'')/u(U) \to u(U')/u(U)$ in $\mathcal{D}/u(U)$
be given. Because $u$ satisfies property (4) there exists
a covering $\{f_i : U_i'' \to U''\}$ in $\mathcal{C}$
and morphisms $c_i : U_i'' \to U'$ such that $u(c_i) = c \circ u(f_i)$.
We think of $U_i''$ as an object over $U$ via the composition
$U_i'' \to U'' \to U$.
It may not be true that $c_i$ is a morphism over $U$!
But since $u(c_i)$ is a morphism over $u(U)$ we may apply
property (3) for $u$ and find coverings $\{f_{ik} : U''_{ik} \to U''_i\}$
such that $c_{ik} = c_i \circ f_{ik} : U''_{ik} \to U'$ are morphisms over $U$.
Hence $\{f_i \circ f_{ik} : U''_{ik}/U \to U''/U\}$ is a covering
in $\mathcal{C}/U$ such that $u(c_{ik}) = c \circ u(f_{ik})$.
Hence (4) holds
for $u : \mathcal{C}/U \to \mathcal{D}/u(U)$.

\medskip\noindent
Property (5). Let $h : V \to u(U)$ be an object of $\mathcal{D}/u(U)$.
Because $u$ satisfies property (5) there exists a covering
$\{c_i : u(U_i) \to V\}$ in $\mathcal{D}$. By property (4) we can find
coverings $\{f_{ij} : U_{ij} \to U_i\}$ and morphisms
$c_{ij} : U_{ij} \to U$ such that $u(c_{ij}) = h \circ c_i \circ u(f_{ij})$.
Hence $\{u(U_{ij})/u(U) \to V/u(U)\}$ is a covering in
$\mathcal{D}/u(U)$ of the desired shape and we conclude that
(5) holds for $u : \mathcal{C}/U \to \mathcal{D}/u(U)$.
\end{proof}

\end{document}
